\scanpage{61}
% @ lines 1–3, row 5 / text
5) $C_{4_{S_0}}:\quad E_0=\mathcal O_{S_0}^3,\quad q_0(x,y,z)=xy+z^2,\quad\omega=e_1\wedge e_2\wedge e_3$ où $(e_1,e_{\unsure{2|3}},e_3)$ est la base canonique (on écrit $(x,y,z)$ pour $xe_1+ye_2+ze_3$)
% @ line 4, row 6 / display
\[
 6)\quad C_{5_{S_0}}:\qquad X_0=\mathbb P^1_{S_0}=\mathbb P(\mathcal O_{S_0}^2)=\check{\mathbb P}(\mathcal O_{S_0}^2)
\]
% @ lines 5–6 / text
$C'_{5_{S_0}}:\quad P_0=\mathbb P^2_{S_0}=\check{\mathbb P}(E_0),\quad C_0=\mathbb V(q_0)$ quadrique définie par la forme $q_0$ de 5°).

% @ line 7, heading / text
II \uline{Opérations de $G_0$}
% @ line 8 / text
1) Translations à gauche
% @ lines 9–11, struck block / deletion
\deleted{\unsure{3|2}) $G_0$ opère sur $\unsure{\mathfrak g|g}_0=\underline{\operatorname{Lie}}(G_0)$ par repr. adjointe. Les op. de $\widetilde G_0$ et $\unsure{\mathscr L_0|L_0|\mathcal L_0}$ sur leurs Algèbres de Lie respectives}
% @ lines 12–15 / text
2) $G_0$ opère sur lui-$=$ par autom. int$^{\mathrm s}$, les op. de $\widetilde G_0$ \deleted{\unsure{et $\mathscr L_0$}} sur \deleted{\unsure{eux|lui}}-mêmes se factorisent par les centres, $Z_0$ \deleted{\unsure{et $G_{m,S_0}$} et} on a des \deleted{isogénies} iso can
% @ line 16 / display
\[
 \widetilde G_0/Z_0=\mathrm{Sl}(2)/\mu_2\simeq\mathrm{GP}(1)=G_0\tag*{(1)}
\]
% @ line 17, crossed formula / display
\[
 \deleted{\mathscr L_0/G_{m,S_0}=\mathrm{Gl}(2)/G_m\simeq\mathrm{GP}(1)=G_0}
\]
% @ line 18 / text
d'où opérations de $G_0$ sur $\widetilde G_0$. \deleted{\unsure{et $\mathscr L_0$}}
% @ lines 19–20 / text
3) Les op. de $G_0$ sur $G_0$ et $\widetilde G_0$ définissent des op. sur $\mathfrak g_0=\underline{\operatorname{Lie}}(G_0)$ et $\widetilde{\mathfrak g}_0=\operatorname{Lie}(\widetilde G_0)$.
% @ lines 21–23 / text
4) L'opération régulière de $\check{\mathbb V}(M_0)^*$ \inserted{$\simeq\mathrm{Gl}(2)$} sur \deleted{\unsure{$M$}} \deleted{\unsure{est}}\inserted{$M_0$} est \deleted{nulle}\inserted{triviale} sur le centre $\simeq G_m\subset\mathrm{Gl}(\unsure{2|1})$, donc passe au quotient en une opération de
% @ line 24 / display
\[
 \mathrm{Gl}(2)/G_m\overset{\mathrm{dfn}}{=}\mathrm{GP}(1)\tag*{(2)}
\]
% @ line 25 / text
sur $M_0$.
% archive sheet number: 60
